\documentclass[12pt,a4paper]{article}

\usepackage[utf8]{inputenc}
\usepackage[T1]{fontenc}
\usepackage[margin=2.5cm]{geometry}
\usepackage{hyperref}
\usepackage{microtype}
\usepackage{parskip}
\usepackage{lmodern}

\hypersetup{colorlinks=true, urlcolor=blue!60!black}

\pagestyle{empty}

\begin{document}

\begin{flushright}
Kundai Farai Sachikonye \\
Auf der Lichtung 19, 82178 Puchheim, Germany \\
\href{https://github.com/fullscreen-triangle}{github.com/fullscreen-triangle} \\
\texttt{kundai.sachikonye@bitspark.com} \\
(+49) 1626386344 \\[6pt]
18 May 2026
\end{flushright}

\vspace{1em}

Prof.\ Dr.\ Astrid Dempfle \\
Institut für Medizinische Statistik \\
Universitätsklinikum Schleswig-Holstein (UKSH) \\
Christian-Albrechts-Universität zu Kiel \\
Kennziffer 28393

\vspace{1.5em}

\textbf{Re: Wissenschaftlicher Mitarbeiter (m/w/d) --- Genetische Epidemiologie /
Statistische Genetik (Kennziffer 28393)}

\vspace{1em}

Dear Prof.\ Dempfle,

I am applying for the position in genetic epidemiology and statistical genetics
at IMS (Kennziffer 28393), in the DFG project on polygenic risk prediction for
individuals with elevated familial IBD risk.
My background spans genome-wide genomic data analysis validated on large
population cohorts, statistical methodology for individual-level disease risk
prediction, and multi-omics integration --- covering the quantitative and
computational requirements the position describes.
I note that the position lists fluent German; mine is conversational rather
than fluent, and I flag this transparently.

\textbf{Genomic data analysis at population scale.}
The \href{https://github.com/fullscreen-triangle/gospel}{Gospel} framework
processes genome-wide variant data at $10^6$ variants per second, integrates
RNA-seq quantification alongside genomic data, and applies Bayesian network
inference and Gaussian mixture models for multi-omics statistical modelling.
Validated against 1000~Genomes Phase~3, gnomAD~v3.1.2, and UK~Biobank;
covers the full pipeline from VCF variant processing through pathway analysis
(KEGG, Reactome) to functional annotation.
The statistical models --- Bayesian networks for genomic variant co-occurrence,
variational Bayes for latent component decomposition, GMMs for population
structure --- are directly applicable to polygenic score construction and
familial risk stratification: they provide a probabilistic framework for
combining many weak genetic signals into a coherent individual-level risk
estimate, which is precisely what the DFG project requires.

\textbf{Individual-level disease risk prediction and structured risk scores.}
Beyond standard polygenic score methodology, I have developed a complementary
approach to individual disease risk prediction in the
\href{https://github.com/fullscreen-triangle/syndrome}{Syndrome} framework.
Disease state is represented as a vector $\mathbf{D}\in[0,1]^8$ across eight
pathophysiological dimensions, each with a per-dimension coherence index
$\eta_i\in[0,1]$ derived from biomarker time-series regularity.
The framework locates the patient's current state in a partition geometry on
the $[0,1]^8$ hypercube and computes an $\mathcal{O}(N\log N)$ trajectory
to the healthy attractor; therapeutic target selection uses sparse $\ell_1$
linear programming to identify the minimum intervention set.
For IBD specifically, the eight dimensions map naturally to the distinct
pathophysiological components of disease risk --- mucosal barrier integrity,
immune cell activation, microbial community structure, stress-response axis,
genetic susceptibility load, epigenetic modification state, cytokine balance,
and circadian regulatory coherence.
A polygenic risk score for IBD captures the genetic susceptibility dimension;
the $\mathbf{D}\in[0,1]^8$ representation integrates it with the remaining
seven dimensions into a single structured risk coordinate.
This is a methodological contribution to the project's aim of risk prediction
in familial settings, where both genetic load and shared environmental factors
are elevated: the framework naturally accommodates non-additive interactions
between genetic and non-genetic risk components that standard linear PRS cannot
represent.
A companion paper proves that loop-holonomy self-consistency is the unique
viable disease diagnostic criterion for this framework, achieving AUC$\,=1.00$
in 25/25 validation experiments.

\textbf{Family data and pedigree-based methods.}
The statistical challenge specific to familial IBD risk prediction is that the
units of analysis are correlated --- family members share genetic background and
often shared environment --- which violates the independence assumptions of
standard GWAS-derived PRS.
My statistical training covers the relevant methodological territory:
mixed-effects models for correlated data, Bayesian hierarchical models for
shared latent structure, and graph-based methods (Kirchhoff circuit laws on
correlation graphs, analogous to kinship matrix eigendecomposition in
linear mixed models) for capturing the dependency structure.
The Gospel framework's Bayesian network layer handles graphical dependency
structures of arbitrary topology, which extends naturally to pedigree graphs.

\textbf{Programming and statistical software.}
I work primarily in Python and R. For statistical genetics specifically,
I have used the standard genomic data formats (VCF, PLINK bed/bim/fam)
and reference databases (gnomAD, 1000~Genomes, UK Biobank) in the Gospel
framework. I am familiar with the PRS methodology landscape (clumping and
thresholding, Bayesian shrinkage approaches such as LDpred, continuous
shrinkage priors) from the literature, though my own implementations use
Bayesian network and variational Bayes approaches rather than PLINK-based
pipelines.
R is my primary language for statistical modelling; Python for data processing
and ML pipeline construction.

\textbf{Background.}
I hold an MSc in Computational Biology (Universität Tübingen) covering
statistical genetics and bioinformatics, an MSc in Pharmaceutical Biotechnology
(HAW Hamburg), and a BSc in Biochemistry and Cell Biology (Jacobs University
Bremen). I conducted doctoral research in Computational Lipidomics at TUM,
covering multi-omics statistical modelling, federated learning for distributed
clinical analysis, and large-scale spectral data pipelines.
I am legally resident in Germany with full work authorisation.
I am available immediately and willing to relocate to Kiel.
German is conversational; I am working to improve it actively.

\vspace{1.5em}
Yours sincerely,

\vspace{2.5em}
\textbf{Kundai Farai Sachikonye}

\end{document}
